\section{The Lee--Gauduchon cone}
\label{sec:lee-gauduchon-cone}

For the rest of the chapter, let $X$ be an arbitrary compact complex manifold of dimension $n\geq2$. To detect pseudoeffectivity in the image of the Lee map $j_X$ of \eqref{eq:lck-lee-bc-map}, one must pair against every Gauduchon metric. The closed forms $\dc(\gamma^{n-1})$ record these tests in de Rham cohomology. Their natural ambient space is the annihilator of $\ker j_X$, where variation of the metric gives an open convex cone and strict duality detects every nonzero positive class in the image.

\Needspace{12\baselineskip}
Set $V_X=H^1_{\mathrm{dR}}(X,\R)$ and
\[
 K_X=\ker j_X,\qquad
 W_X=K_X^\perp\subset H^{2n-1}_{\mathrm{dR}}(X,\R),
\]
where the annihilator uses the Poincar\'e pairing with a class $a\in H^{2n-1}_{\mathrm{dR}}(X,\R)$
\[
 \langle[\lambda],a\rangle=\int_X\lambda\wedge a.
\]
Finally, let
\[
 \mathcal C_X
 =\{[\lambda]\in V_X:j_X([\lambda])\text{ is pseudoeffective}\}.
\]
The image of $j_X$ can be strictly smaller than the kernel of the Bott--Chern-to-de-Rham map, as \cref{ex:lee-calabi-eckmann} will show.

\begin{lemma}\label{lem:lee-gauduchon-kernel}
The space $K_X$ consists of the de Rham classes of real parts of closed holomorphic one-forms. Moreover,
\[
 \mathcal C_X\cap(-\mathcal C_X)=K_X.
\]
\end{lemma}
\begin{proof}
If $j_X([\lambda])=0$, write $\dc\lambda=\ddc f$ for a smooth real function $f$. Since $\dc\mathrm{d}=-\mathrm{d}\dc$, the form $\lambda+\mathrm{d}f$ is both $\mathrm{d}$-closed and $\dc$-closed. Its $(1,0)$-part is therefore holomorphic and closed: the two equations separate the $\partial$ and $\bar\partial$ derivatives. Thus $\lambda+\mathrm{d}f$ is the real part of the closed holomorphic form $2(\lambda+\mathrm{d}f)^{1,0}$. Conversely, the real part of a closed holomorphic form is closed for both operators, and its image under $j_X$ is zero.

By \cref{prop:psef-cone-compact-base}, the pseudoeffective Bott--Chern cone contains no nonzero line. Thus $j_X(\xi)$ and $-j_X(\xi)$ are both pseudoeffective exactly when $j_X(\xi)=0$. This is the asserted identity $\mathcal C_X\cap(-\mathcal C_X)=K_X$.
\end{proof}

\Needspace{5\baselineskip}
\begin{definition}\label{def:lee-gauduchon-cone}\index{Lee--Gauduchon cone}
The \emph{Lee--Gauduchon cone} is
\[
 \operatorname{LG}(X)
 =\bigl\{[\dc(\gamma^{n-1})]:
                    \gamma\text{ is a Gauduchon metric on }X\bigr\}.
\]
\end{definition}

For a closed one-form $\lambda$ and any smooth real $(n-1,n-1)$-form $\Gamma$ with $\ddc\Gamma=0$, integration by parts gives
\begin{equation}\label{eq:lck-lee-gauduchon-pairing}
 \int_X\lambda\wedge\dc\Gamma
 =\int_X\kappa_\lambda\wedge\Gamma.
\end{equation}
The form $\dc\Gamma$ is closed, and both sides depend only on the relevant cohomology classes. In particular the class on the left annihilates $K_X$.

\begin{proposition}\label{prop:lee-gauduchon-open}
The cone $\operatorname{LG}(X)$ is nonempty, convex and open in $W_X$. Its linear span is $W_X$.
\end{proposition}
\begin{proof}
Let $\mathcal Z_X$ be the real vector space of smooth $(n-1,n-1)$-forms $\Gamma$ satisfying $\ddc\Gamma=0$, and set
\[
 q_X(\Gamma)=[\dc\Gamma].
\]
First we show $q_X(\mathcal Z_X)=W_X$. The inclusion into $W_X$ follows from \eqref{eq:lck-lee-gauduchon-pairing}. If $[\lambda]$ annihilates the image of $q_X$, its pairing with every Gauduchon form is zero. The generalized Lamari criterion, \cref{thm:lamari}, implies that both $\kappa_\lambda$ and $-\kappa_\lambda$ have pseudoeffective Bott--Chern classes. By \cref{lem:lee-gauduchon-kernel}, $[\lambda]\in K_X$. Thus the image and $W_X$ have the same annihilator in $V_X$. The finite-dimensional nondegenerate Poincar\'e pairing proves their equality.

Let $\mathcal Z_X^+$ denote the strictly positive forms in $\mathcal Z_X$. Every such form has a unique positive $(n-1)$st root, by the matrix construction used in the proof of \cref{thm:lamari}. This root is Gauduchon. Hence
\[
 \operatorname{LG}(X)=q_X(\mathcal Z_X^+).
\]
Gauduchon's theorem makes $\mathcal Z_X^+$ nonempty. It is convex and stable under positive scaling.

To see openness directly, choose forms $\Gamma_1,\ldots,\Gamma_r\in\mathcal Z_X$ whose images form a basis of $W_X$. If $\Gamma\in\mathcal Z_X^+$, then $\Gamma+\sum_i t_i\Gamma_i$ remains strictly positive for all sufficiently small real $t_i$. Its image describes a neighborhood of $q_X(\Gamma)$ in $W_X$. This also covers $W_X=\{0\}$, with the empty basis. Since a nonempty open subset spans its ambient vector space, the span assertion follows.
\end{proof}

\Needspace{12\baselineskip}
\begin{theorem}\label{thm:lee-gauduchon-duality}\index{Lee--Gauduchon cone!duality}
The cone $\mathcal C_X$ is closed, and
\[
 \mathcal C_X
 =\{\xi\in V_X:\langle\xi,a\rangle\geq0
                         \text{ for every }a\in\operatorname{LG}(X)\}.
\]
Conversely,
\[
 \operatorname{LG}(X)
 =\{a\in W_X:\langle\xi,a\rangle>0
                         \text{ for every }\xi\in\mathcal C_X\setminus K_X\}.
\]
\end{theorem}
\begin{proof}
The first equality is exactly \cref{thm:lamari} applied to the smooth closed background $\kappa_\lambda$, using \eqref{eq:lck-lee-gauduchon-pairing}. It expresses $\mathcal C_X$ as an intersection of closed half-spaces, so this cone is closed.

If $a=[\dc(\gamma^{n-1})]\in\operatorname{LG}(X)$ and $\xi=[\lambda]\in\mathcal C_X\setminus K_X$, choose $T=\kappa_\lambda+\ddc u\geq0$. This current is nonzero because it represents the nonzero class $j_X(\xi)$. Therefore
\[
 \langle\xi,a\rangle
 =\int_XT\wedge\gamma^{n-1}>0.
\]

For the converse, suppose $a\in W_X\setminus\operatorname{LG}(X)$. Finite-dimensional separation of a point from a nonempty open convex cone gives a nonzero functional $\ell\in W_X^*$ such that
\[
 \ell(a)\leq0,\qquad \ell(b)\geq0
                   \quad(b\in\operatorname{LG}(X)).
\]
Indeed, in the usual separating inequalities the constant can be taken to be zero by scaling elements of the cone. By Poincar\'e duality, $\ell$ is represented by a class $\xi\in V_X$, modulo $K_X$. Its nonzero restriction means $\xi\notin K_X$. The first equality in the theorem implies $\xi\in\mathcal C_X$, contradicting the required strict inequalities for $a$. This proves the second equality. If $\operatorname{LG}(X)=W_X$, there is no point to separate and the same conclusion holds.
\end{proof}

\subsection{Bimeromorphic invariance}

A modification preserves degree-one cohomology and detects pseudoeffectivity of classes pulled back from the base. These two facts identify the cones $\mathcal C_X$; duality then identifies the Lee--Gauduchon cones.

\Needspace{10\baselineskip}
\begin{lemma}\label{lem:lee-gauduchon-modification}
Let $\pi:Y\to X$ be a modification between connected compact complex manifolds, and let $\theta$ be a smooth closed real $(1,1)$-form on $X$. Then
\[
 \{\theta\}\text{ is pseudoeffective}
 \quad\Longleftrightarrow\quad
 \{\pi^*\theta\}\text{ is pseudoeffective}.
\]
The pullback on real Bott--Chern cohomology in bidegree $(1,1)$ is injective.
\end{lemma}
\begin{proof}
By \cref{prop:psef-cone-compact-base}, both pullback and pushforward preserve pseudoeffectivity. The underlying potential pullback is the one used in \cref{lem:modification-lone-pullback}. Since $\pi$ has degree one,
\[
 \pi_*\pi^*\{\theta\}=\{\theta\}.
\]
Indeed, the current identity $\pi_*\pi^*\theta=\theta$ follows by integration on the isomorphism locus, whose complement has smooth volume zero. This left inverse proves injectivity and the reverse implication for pseudoeffectivity.
\end{proof}

The degree-one comparison follows by moving cycles off the exceptional image.

\begin{lemma}\label{lem:lee-gauduchon-blowup-cohomology}
If $\pi:Y\to X$ is a modification between connected compact complex manifolds, then
\[
 \pi^*:H^1_{\mathrm{dR}}(X,\R)\longrightarrow
                    H^1_{\mathrm{dR}}(Y,\R)
\]
is an isomorphism. Its Poincar\'e dual is the isomorphism
$\pi_*:H^{2n-1}_{\mathrm{dR}}(Y,\R)\to H^{2n-1}_{\mathrm{dR}}(X,\R)$.
\end{lemma}
\begin{proof}
For a proper analytic subset $A$ of a connected complex manifold $M$, restriction
$H^1(M,\R)\to H^1(M\setminus A,\R)$ is injective, and is an isomorphism
if $A$ has complex codimension at least two. Indeed, choose a Whitney
stratification of $A$. General position moves every one-cycle off its
strata, which have real codimension at least two. If those codimensions
are at least four, it also moves every bounding two-chain off the strata,
relative to its boundary. Thus the map on first homology is surjective
in the first case and bijective in the second. Dualizing over $\R$ gives
the asserted restriction statements.

Since $X$ is smooth, there is an analytic subset $A\subset X$ of complex
codimension at least two such that $\pi$ is an isomorphism over
$U=X\setminus A$. Restriction from $X$ to $U$ is therefore an isomorphism,
whereas restriction from $Y$ to $\pi^{-1}U$ is injective. Every class on
$\pi^{-1}U\simeq U$ extends to $X$ and hence pulls back to a class on $Y$.
The commutative restriction diagram proves that $\pi^*$ is an isomorphism.
Poincar\'e duality and the projection formula give the assertion for $\pi_*$.
\end{proof}

\Needspace{10\baselineskip}
\begin{theorem}\label{thm:lee-gauduchon-bimeromorphic}
Let $X$ and $Y$ be bimeromorphic connected compact complex manifolds. Under the natural identification of $H^1_{\mathrm{dR}}$ and its Poincar\'e-dual identification of $H^{2n-1}_{\mathrm{dR}}$, their spaces $K_X,W_X$ and cones $\mathcal C_X,\operatorname{LG}(X)$ agree.
\end{theorem}
\begin{proof}
First let $\pi:Y\to X$ be a modification. Holomorphic pullback commutes with $\dc$, so
\[
 j_Y(\pi^*\xi)=\pi^*j_X(\xi).
\]
By \cref{lem:lee-gauduchon-modification}, membership in $\mathcal C_X$ is equivalent to membership of the pullback in $\mathcal C_Y$, and membership in $K_X$ is likewise preserved in both directions. By \cref{lem:lee-gauduchon-blowup-cohomology}, every degree-one class upstairs occurs in this way. Thus the degree-one isomorphism identifies the two cones and kernels. Its dual identifies $W_Y$ with $W_X$, and \cref{thm:lee-gauduchon-duality} gives
\[
 \pi_*\operatorname{LG}(Y)=\operatorname{LG}(X).
\]

For a bimeromorphic map between $X$ and $Y$, resolve its graph to obtain
modifications $p:Z\to X$ and $q:Z\to Y$ from one smooth compact manifold.
The preceding argument applies to both maps and identifies the cones and
kernels via $(q^*)^{-1}p^*$ and its Poincar\'e dual. These identifications
are independent of the resolution: they agree on a common open set where
the original map is an isomorphism, and the injectivity of degree-one
restriction proved above determines each class uniquely. The dual
identifications are consequently canonical as well.
\end{proof}

\subsection{Balanced metrics and the Hopf cone}

The presence of zero in the Lee--Gauduchon cone rules out nonzero pseudoeffective classes in the image of $j_X$. Balanced metrics satisfy this condition. On the Hopf threefold, by contrast, the Lee curvature gives a positive ray and the dual cone is an open half-line.

\begin{corollary}\label{cor:lee-gauduchon-zero}
The following conditions are equivalent:
\[
 0\in\operatorname{LG}(X),\qquad
 \operatorname{LG}(X)=W_X,\qquad
 \mathcal C_X=K_X.
\]
A balanced metric implies these conditions. On a manifold in Fujiki's class $\mathcal C$, one has $W_X=\{0\}$ and $\operatorname{LG}(X)=\{0\}$.
\end{corollary}
\begin{proof}
An open cone in a finite-dimensional vector space that contains zero is the entire space: scale a neighborhood of zero to reach any point. The two dual descriptions in \cref{thm:lee-gauduchon-duality} now prove the equivalences, including the case $W_X=\{0\}$. If $b$ is balanced, its form $b^{n-1}$ is closed and of pure type, hence also $\dc$-closed. It is Gauduchon and represents zero in $\operatorname{LG}(X)$.

On a compact K\"ahler manifold the $\partial\bar\partial$-lemma gives $j_X=0$: every $\kappa_\lambda$ is closed and de Rham exact. Thus $K_X=V_X$ and $W_X=0$. Nonemptiness gives $\operatorname{LG}(X)=\{0\}$. Bimeromorphic invariance extends this conclusion to Fujiki's class $\mathcal C$.
\end{proof}

Consequently, a nonzero pseudoeffective class in the image of $j_X$ obstructs balanced metrics throughout the bimeromorphic class. For LCK manifolds, \cref{conj:lck-lee-psef} would provide such a class.

For the Hopf threefold of \cref{chap:hopf}, the covering description identifies the underlying manifold with $S^1\times S^5$, so $V_X=\R[\lambda]$. The class $j_X([\lambda])=\{\alpha\}$ is nonzero by \eqref{eq:lck-lee-degree}, and is pseudoeffective because $\alpha\geq0$. Negative multiples are not pseudoeffective, since their pairing against the fixed Gauduchon metric is negative. Thus
\[
 K_X=0,\qquad \mathcal C_X=\R_{\geq0}[\lambda],\qquad
 \operatorname{LG}(X)=\{a\in H^5_{\mathrm{dR}}(X,\R):
                                  \langle[\lambda],a\rangle>0\}.
\]
The same manifold has unbounded Monge--Amp\`ere mass by \cref{thm:hopf-unbounded}. The positive ray detected here controls linear cohomological pairings, which do not bound those nonlinear masses.
